%! Sinusoidal Functions — Unit 3, Lesson 3.5 — Exit Ticket
\documentclass[10pt]{article}
\usepackage{precalculus-article}
\usepackage{precalculus-boxes}
\newcommand{\TallMath}[1]{$\displaystyle #1\rule[-1.4em]{0pt}{3.2em}$}

\begin{document}

\pageheader{Unit 3, Lesson 3.5}{Exit Ticket}
\namedateperiod

\vspace{0.18in}

\begin{notesbox}{Exit Ticket: Sinusoidal Functions}

\textbf{1.} A sinusoidal function has a maximum value of $7$ and a minimum value
of $1$.

\medskip
(a) What is the amplitude of this function?

\vspace{0.08in}
\begin{tabular}{@{}l@{}}
  Amplitude $= \dfrac{\blank{1.5cm} - \blank{1.5cm}}{2} = $ \blank{2cm}
\end{tabular}

\vspace{0.18in}

(b) What is the midline of this function?  Write your answer as an equation.

\vspace{0.08in}
\begin{tabular}{@{}l@{}}
  Midline: $y = \dfrac{\blank{1.5cm} + \blank{1.5cm}}{2} = $ \blank{2cm}
\end{tabular}

\vspace{0.25in}

\textbf{2.} The period of a sinusoidal function is $\dfrac{\pi}{2}$.
What is the frequency of this function?

\vspace{0.1in}
\begin{tabular}{@{}l@{}}
  Frequency $= \dfrac{1}{\blank{2cm}} = $ \blank{3cm}
\end{tabular}

\end{notesbox}

\end{document}
